% ==============================================================================
% ANNEXE 1 -- GUIDE D'ENTRETIEN SEMI-DIRECTIF
% Mémoire MAE -- Arthur SAUNIER -- Wavestone Luxembourg
% Le \chapter{} est déclaré dans le fichier maître ; ce fichier démarre en \section.
% ==============================================================================

\section{Objectifs et architecture du guide}

Le présent guide constitue l'instrument de collecte de la phase empirique. Il a été
construit de manière déductive à partir des quatre questions directrices dérivées de
la question centrale de recherche, chaque thème d'entretien constituant la
traduction conversationnelle d'une de ces questions. Ce chaînage explicite entre
questionnement théorique et questionnement de terrain garantit la validité interne
du dispositif~: aucune question posée aux répondants n'est orpheline d'une
interrogation de recherche, et aucune question de recherche n'est laissée sans
matériau empirique.

L'entretien est calibré pour une durée de trente-cinq à quarante-cinq minutes,
compatible avec les contraintes de disponibilité de consultants en mission. Cette
contrainte impose une discipline de conduite~: les thèmes~3 et~4 étant les plus
exposés au risque d'être sacrifiés en fin d'entretien, l'enquêteur veille à borner
strictement les deux premiers.

\begin{table}[h]
\centering
\small
\begin{tabular}{|p{0.08\linewidth}|p{0.40\linewidth}|p{0.23\linewidth}|p{0.09\linewidth}|}
\hline
\textbf{Thème} & \textbf{Question directrice adossée} & \textbf{Cadre théorique mobilisé} & \textbf{Durée} \\
\hline
T0 & Cadrage du répondant (hors questionnement) & --- & 5 min \\
\hline
T1 & QD1 -- Appropriation réelle des outils d'IA générative dans les tâches d'analyse & Sociologie des usages, TAM (Davis, 1989) & 10 min \\
\hline
T2 & QD2 -- Impact sur l'apprentissage par la pratique et la montée en compétences & Nonaka \& Takeuchi (1995), Polanyi (1966) & 10 min \\
\hline
T3 & QD3 -- Reconstruction de la légitimité de l'expert face au client & Abbott (1988), Suchman (1995) & 8 min \\
\hline
T4 & QD4 -- Encadrement managérial des risques sans bridage de l'autonomie & Mintzberg (1982), gouvernance des risques & 8 min \\
\hline
T5 & Clôture projective et ouverture & --- & 4 min \\
\hline
\end{tabular}
\caption{Architecture du guide d'entretien et adossement théorique}
\label{tab:architecture-guide}
\end{table}

\section{Consignes de conduite adressées à l'enquêteur}

La qualité du matériau recueilli dépend moins de la formulation des questions que de
la discipline de l'enquêteur. Les principes suivants encadrent la conduite de chaque
entretien.

\begin{itemize}
    \item \textbf{Privilégier le récit à l'opinion.} Chaque thème s'ouvre par une
    question de type narratif ou par un incident critique (« racontez-moi la
    dernière fois que… ») plutôt que par une question d'opinion. Le récit d'une
    situation vécue produit un matériau vérifiable et échappe aux réponses de
    façade que suscite un sujet socialement connoté.

    \item \textbf{Ne jamais introduire le vocabulaire de la recherche.} Les termes
    « légitimité », « montée en compétences », « hallucination », « gouvernance »
    ne doivent pas être prononcés par l'enquêteur avant que le répondant ne les ait
    lui-même mobilisés. Leur emploi prématuré induit la réponse et détruit la valeur
    probante du verbatim.

    \item \textbf{Tenir une position de neutralité axiologique.} L'enquêteur, lui-même
    consultant du cabinet, s'abstient de défendre comme de critiquer l'usage de ces
    outils. Toute marque d'approbation ou de réprobation oriente durablement la suite
    de l'entretien.

    \item \textbf{Exploiter le silence.} Une pause de trois à cinq secondes après une
    réponse jugée terminée produit fréquemment la relance la plus riche de l'entretien.

    \item \textbf{Relancer par la reformulation.} Les relances privilégiées sont
    l'écho (reprise du dernier segment de phrase sur un ton interrogatif), la demande
    de précision factuelle (« sur quelle mission~? ») et la demande d'exemple
    contraire (« et une fois où cela ne s'est pas passé ainsi~? »).

    \item \textbf{Assumer la souplesse de l'ordre.} L'ordre des thèmes est indicatif.
    Si le répondant aborde spontanément un thème ultérieur, l'enquêteur le suit et
    coche mentalement le thème traité plutôt que de le ramener au fil prévu.

    \item \textbf{Ne pas combler les contradictions.} Les incohérences entre le
    discours déclaré et les pratiques décrites constituent le matériau analytique le
    plus précieux~; elles sont notées, non résolues en séance.
\end{itemize}

\section{Préambule, anonymisation et recueil du consentement}

Le texte suivant est lu au répondant avant tout enregistrement.

\begin{quote}
\textit{Merci d'avoir accepté cet entretien. Il s'inscrit dans le cadre d'un mémoire
de Master~2 en Management et Administration des Entreprises, réalisé à l'IAE Nancy
parallèlement au stage effectué dans le cabinet. Ce travail porte sur la manière dont
les consultants intègrent les outils d'intelligence artificielle générative dans leur
travail d'analyse, et sur ce que cela change dans l'exercice du métier.}

\textit{Trois précisions avant de commencer. Premièrement, il ne s'agit ni d'un audit,
ni d'une évaluation~: il n'existe pas de bonne ou de mauvaise réponse, et la
description de pratiques informelles est aussi utile que celle des usages officiels.
Deuxièmement, l'entretien est intégralement anonymisé~: aucun nom de personne, de
client ou de mission ne figurera dans le mémoire, et les répondants y sont désignés
par un code de type « consultant junior~1 ». Troisièmement, l'entretien dure environ
quarante minutes et l'enregistrement audio, s'il est autorisé, sert uniquement à la
retranscription et est détruit après l'analyse.}

\textit{Acceptez-vous que l'entretien soit enregistré~?}
\end{quote}

En cas de refus d'enregistrement, l'entretien se poursuit sous prise de notes
manuscrite, la restitution des verbatims étant alors signalée comme reconstituée.

\section{Thème 0 --- Trajectoire et cadre d'exercice}

Ce thème ne relève d'aucune question directrice. Il vise à installer la parole, à
recueillir les variables de contrôle nécessaires à la lecture croisée des profils et à
identifier le vocabulaire propre au répondant, que l'enquêteur réemploiera ensuite.

\begin{itemize}
    \item Pour commencer, pouvez-vous décrire votre parcours et la manière dont vous
    êtes arrivé dans la practice~?
    \item Quel type de missions occupe l'essentiel de votre temps en ce moment~?
    \item Sur une mission typique, en quoi consiste concrètement votre journée~?
    \item Quelle part de votre travail relève de la production de livrables écrits~?
\end{itemize}

\section{Thème 1 --- L'usage réel des outils dans le travail d'analyse}

\textit{Question directrice~1~: de quelle manière les consultants s'approprient-ils
les outils d'intelligence artificielle générative dans la réalisation de leurs tâches
d'analyse quotidiennes~?}

\subsection*{Question d'ouverture}

Pouvez-vous me raconter la dernière fois que vous avez utilisé un de ces outils dans
votre travail~? Que cherchiez-vous à faire, et comment cela s'est-il passé~?

\subsection*{Relances de premier niveau}

\begin{itemize}
    \item Sur quelles tâches précisément ces outils interviennent-ils dans votre
    quotidien~?
    \item Comment avez-vous appris à vous en servir~? Par une formation, par un
    collègue, seul~?
    \item À quel moment de la production d'un livrable les mobilisez-vous~: au
    démarrage, en cours de rédaction, à la relecture~?
    \item Y a-t-il des tâches pour lesquelles vous ne les utilisez jamais~? Lesquelles,
    et pourquoi~?
    \item Utilisez-vous les outils mis à disposition par le cabinet, ou d'autres~?
    Qu'est-ce qui fait la différence pour vous~?
\end{itemize}

\subsection*{Relances de second niveau}

À n'utiliser que si le répondant demeure dans la généralité.

\begin{itemize}
    \item Vous parlez de gain de temps~: sur la dernière mission, à quoi avez-vous
    consacré le temps ainsi libéré~?
    \item Y a-t-il des situations où cela vous a demandé plus d'effort que de le faire
    vous-même~?
    \item Vous êtes-vous déjà arrêté en cours d'utilisation en vous disant que ce
    n'était pas la bonne méthode~? Qu'est-ce qui a déclenché ce moment~?
    \item En parlez-vous avec vos collègues~? Est-ce un sujet dont on discute
    ouvertement dans l'équipe~?
\end{itemize}

\section{Thème 2 --- Apprentissage du métier et construction de l'expertise}

\textit{Question directrice~2~: quel est l'impact de cette automatisation sur
l'apprentissage par la pratique et la trajectoire de montée en compétences des
consultants juniors~?}

\subsection*{Question d'ouverture}

Si vous repensez à la manière dont vous avez appris ce métier --- ou dont vous voyez
les nouveaux arrivants l'apprendre aujourd'hui ---, qu'est-ce qui a changé~?

\subsection*{Relances de premier niveau}

\begin{itemize}
    \item Quelles sont, selon vous, les tâches par lesquelles on apprend réellement le
    métier de consultant~?
    \item Ces tâches sont-elles encore réalisées de la même manière aujourd'hui~?
    \item Comment vous y prenez-vous pour vérifier ce que produit l'outil~? Que faut-il
    savoir pour être capable de le vérifier~?
    \item Un consultant qui débuterait aujourd'hui avec ces outils apprendrait-il les
    mêmes choses que vous~?
    \item Comment se passe la relecture d'un livrable par un senior~: qu'est-ce qui se
    transmet à ce moment-là~?
\end{itemize}

\subsection*{Relances de second niveau}

\begin{itemize}
    \item Vous êtes-vous déjà retrouvé à valider un contenu que vous n'auriez pas su
    produire vous-même~? Comment avez-vous procédé~?
    \item Y a-t-il des compétences que vous sentez que vous exercez moins qu'avant~?
    \item Si l'accès à ces outils était coupé demain sur une mission, qu'est-ce qui
    deviendrait difficile pour l'équipe~?
\end{itemize}

\section{Thème 3 --- La relation au client et la valeur de l'expertise}

\textit{Question directrice~3~: comment l'autorité et la légitimité de l'expert humain
se reconstruisent-elles vis-à-vis du client lorsque la machine produit une partie
substantielle de la valeur intellectuelle~?}

\subsection*{Question d'ouverture}

Quand vous présentez un livrable à un client, sur quoi repose selon vous sa confiance
dans ce que vous lui remettez~?

\subsection*{Relances de premier niveau}

\begin{itemize}
    \item Le sujet de l'usage de ces outils a-t-il déjà été abordé avec un client~?
    Dans quelles circonstances, et comment cela s'est-il passé~?
    \item Diriez-vous à un client qu'une partie de l'analyse a été produite avec ces
    outils~? Qu'est-ce qui guiderait cette décision~?
    \item Est-ce que les clients eux-mêmes utilisent ces outils~? Cela change-t-il
    quelque chose à ce qu'ils attendent du cabinet~?
    \item Qu'est-ce que le cabinet vend aujourd'hui que le client ne pourrait pas
    obtenir autrement~?
\end{itemize}

\subsection*{Relances de second niveau}

\begin{itemize}
    \item Avez-vous déjà eu le sentiment qu'un livrable était perçu comme moins
    crédible~? Qu'est-ce qui l'a provoqué~?
    \item Dans dix ans, sur quoi un client acceptera-t-il encore de payer un cabinet
    de conseil~?
\end{itemize}

\section{Thème 4 --- Encadrement, contrôle qualité et gestion des risques}

\textit{Question directrice~4~: comment le management encadre-t-il l'usage de l'IA
pour pallier les risques opérationnels sans brider l'autonomie des collaborateurs~?}

\subsection*{Question d'ouverture}

Existe-t-il des règles, écrites ou non, sur l'usage de ces outils dans le cabinet~?
Comment en avez-vous eu connaissance~?

\subsection*{Relances de premier niveau}

\begin{itemize}
    \item Que se passe-t-il concrètement avant qu'un livrable ne parte chez le client~?
    Qui contrôle quoi~?
    \item Avez-vous déjà rencontré un contenu produit par un outil qui s'est révélé
    faux~? Que s'est-il passé ensuite~?
    \item Comment traitez-vous la question des données de mission lorsque vous
    utilisez ces outils~?
    \item Ces règles vous paraissent-elles adaptées au travail réel~? Y a-t-il des
    situations où elles sont contournées~?
    \item Selon vous, qui est responsable si une erreur passe dans un livrable remis
    au client~?
\end{itemize}

\subsection*{Relances de second niveau}

\begin{itemize}
    \item Si vous aviez à écrire vous-même la règle, que diriez-vous~?
    \item Y a-t-il un équilibre à trouver entre encadrer et laisser faire~? Où le
    placeriez-vous~?
\end{itemize}

\section{Thème 5 --- Clôture}

\begin{itemize}
    \item Si vous étiez à la place de la direction du cabinet, quelle serait la
    première décision que vous prendriez sur ce sujet~?
    \item Y a-t-il un aspect important que nous n'avons pas abordé et que vous
    souhaiteriez ajouter~?
    \item Voyez-vous une personne du cabinet dont le point de vue serait
    particulièrement utile à recueillir~?
\end{itemize}

L'entretien se clôt par le rappel des modalités d'anonymisation et par la proposition
d'une transmission de la synthèse finale au répondant.

\section{Variantes de formulation selon le profil}

Le guide demeure identique dans sa structure quel que soit le répondant, condition de
la comparabilité des matériaux. Seules les formulations d'ancrage sont adaptées afin
de solliciter la position d'énonciation propre à chaque profil.

Le panel distingue deux positions, et cette distinction épouse celle du dispositif de
revue décrit au chapitre~1. Les \textbf{juniors} — analystes et consultants — rédigent les
livrables. Les \textbf{encadrants} — consultants seniors, manager et senior manager —
accompagnent la production puis la contrôlent avant l'envoi au client. Le contraste
recherché est celui de deux positions dans la chaîne de production, non celui de deux
échelons hiérarchiques.

\begin{table}[h]
\centering
\small
\begin{tabular}{|p{0.08\linewidth}|p{0.38\linewidth}|p{0.38\linewidth}|}
\hline
\textbf{Thème} & \textbf{Juniors (analystes, consultants)} & \textbf{Encadrants (consultants seniors, manager, senior manager)} \\
\hline
T1 & Usage vécu et découverte des outils & Usage propre \emph{et} usages observés dans l'équipe \\
\hline
T2 & Sentiment sur son propre apprentissage & Perception de l'évolution des juniors encadrés \\
\hline
T3 & Position en soutenance de livrable & Position en relation client directe \\
\hline
T4 & Connaissance et vécu des règles & Rôle effectif dans le contrôle qualité, arbitrage entre risque et productivité \\
\hline
\end{tabular}
\caption{Adaptation des angles d'interrogation selon la position du répondant}
\label{tab:variantes-profils}
\end{table}

\section{Fiche de traçabilité de l'entretien}

Une fiche est renseignée immédiatement après chaque entretien, avant toute
retranscription, afin de conserver les éléments contextuels non enregistrés.

\begin{itemize}
    \item Code du répondant, profil et ancienneté approximative dans le conseil.
    \item Date, durée effective, modalité (présentiel ou visioconférence).
    \item Autorisation d'enregistrement accordée ou refusée.
    \item Thèmes effectivement couverts et thèmes écourtés.
    \item Climat de l'entretien et éventuelles réticences perçues.
    \item Trois formulations marquantes notées à chaud.
    \item Hypothèses émergentes à tester lors des entretiens suivants.
\end{itemize}